\section{Ran upgrades}
\label{section:Ran_upgrades}

\subsubsection{} In this section we will upgrade our previous constructions to the Ran setting.

\subsection{The Ran version of $\eta_G$}

\subsubsection{} For convenience, we will write
  \[\Poinc_? = \PsId_! \cdot \coeff_*^\vee \cdot \PsId_\Whit^{-1}.\]
We write $\Poinc_{\Ran,?}$ and $\Poinc_{I,?}$ the variants of this functor over $\Ran$ and $X^I$.

\subsubsection{} Consider the diagram
\[
\begin{tikzcd}[
  column sep={8em,between origins},
  row sep={4em,between origins}
]
  \Whit(G)_\Ran^\vee \arrow[r] \arrow[d, "\coeff_*^\vee" left] & \Whit(G)_\Ran \arrow[d, "\Poinc_!"] \\
  \DMod(\Bun_G)^\vee \arrow[r, "\PsId_!"] & \DMod(\Bun_G)
\end{tikzcd}
\]
of $\Sph_{G,\Ran}$-linear categories. The goal of this section is to extend $\eta_G$ to a map
  \[\eta_{G,\Ran} : \Poinc_{\Ran,!} \to \Poinc_{\Ran, ?} [2\delta_G]\]
of $\Sph_{G,\Ran}$-linear functors. We will write the construction over each power of the curve.

\subsubsection{} In fact we can simply repeat the definition of $\eta_G$ in order to define a map
  \[\eta_{G,I} : \Poinc_{I,!} \to \Poinc_{I,?} [2\delta_G]\]
of plain functors. That is, we consider the diagram
\[
\begin{tikzcd}[
  column sep={12em,between origins},
  row sep={4em,between origins}
]
  \left( \Bunc_{N,I}^\pol \underset{\Bun_G}{\times} \Bunc_G \right)^\trans \arrow[r, hook] & \Bunc_{N,I}^\pol \underset{\Bun_G}{\times} \Bunc_G \arrow[d, "\pi_2"] \\
   & \Bun_G
\end{tikzcd}
\]
and define $\eta_{G,I}$ to be
\[
\begin{tikzcd}[
  column sep={6em,between origins},
  row sep={4em,between origins}
]
  & \pi_{2, !} \cdot \left( \WW \stoximes r_{G,!} \right)^\trans \arrow[dr] \arrow[dl, "\sim"] & & \pi_{2, *} \cdot \left( \WW \shoximes r_{G,!} \right) \arrow[dl, "\sim"] \\
  \pi_{2, !} \cdot \left( \WW \stoximes r_{G,!} \right) [2\delta_G] & & \pi_{2, *} \cdot \left( \WW \shoximes r_{G,!} \right)^\trans [2\delta_G]&
\end{tikzcd}
\]
The left arrow is obviously an isomorphism, and the right arrow is an isomorphism by the same argument as Proposition \ref{proposition:cis_vanishing}.

\begin{proposition}
\label{proposition:Sph_linear}
The map $\eta_{G,I}$ is canonically $\Sph_{G,I}$-linear.
\end{proposition}

\subsubsection{Remarks}
\begin{itemize}
    \item[(i)] This result is not tautological. It requires the functorial construction of commutative diagrams:
\[
\begin{tikzcd}[
  column sep={10em,between origins},
  row sep={4em,between origins}
]
  \SS \star ( \Poinc_! \cdot \WW) \arrow[r] \arrow[d, "\sim"] & \SS \star (\Poinc_? \cdot \WW) \arrow[d, "\sim"] \\
  \Poinc_! \cdot (\SS \star \WW) \arrow[r] & \Poinc_? \cdot (\SS \star \WW)
\end{tikzcd}
\]
The right vertical arrow uses the fact that $\PsId_{G,I,!}$ is linear for the action of $\Sph_{G,I}$. It is easy to prove this by observing that 
\[
    \PsId_{G,I,!} \star \SS \simeq \SS \star \PsId_{G,I,!} \in \DMod(\Bun_G \times \Bun_G \times X^I)
\]
is given by the $!$-pushforward of $\SS$ along the map
  \[\hl \times \hr \times s : \HH_{G,I} \to \Bun_G \times \Bun_G \times X^I.\]
But this identification is difficult to use because has the wrong (i.e. $*$-pull and $!$-push) nature. Therefore, we should remedy this by constructing a compactification of $\hl \times \hr \times s$ as an $(\HH_{G, I}, \HH_{G,I})$-bimodule.
    \item[(ii)] Since $\Whit(G)_\Ran$ is freely generated by $\WW_\vac$ under the action of $\Rep(\LG)_\Ran$, we automatically obtain a $\Rep(\LG)_\Ran$-linear version of $\eta_{G,\Ran}$. But this map is not obviously $\Sph_{G, \Ran}$-linear. Furthermore, it is not clear that this na\"ively defined map agrees with $\eta_{G,\Ran}$ as a plain natural transformation and this identification will be needed in the sequel.
\end{itemize}

\subsubsection{} We write
  \[\Vin\Bun_{G,I}^\pol \to \Bun_G \times \Bun_G \times X^I\]
for the stack classifying a point $x \in X^I$ along with the following mapping data:
\[
\begin{tikzcd}[
  column sep={8em,between origins},
  row sep={4em,between origins}
]
  \eta \arrow[r, hook] \arrow[d] & X \setminus x \arrow[r, hook] \arrow[d] & X \arrow[d] \\ 
  G \bs \Vin^\circ_G / G \arrow[r, hook] & G \bs \Vin_G / G \arrow[r] & G \bs T_\ad^+ / G
\end{tikzcd}
\]
We remind the reader that the action of $G \times G$ on $T_\ad^+$ appearing the bottom right of this diagram is trivial. As with $\Vin\Bun_G$, this stack comes with an action of $T$ and we define
  \[\Bunc_{G,I}^\pol = \Vin\Bun_{G,I}^\pol / T.\]
Here are the salient geometric properties of this stack:

\begin{proposition}
\label{proposition:Hecke_comp_propoerties}
We have the following:
\begin{itemize}
    \item[(i)] There is a factorization
      \[\Bunc_{G,I} \mon \Bunc_{G,I}^\pol \to \Bun_G \times \Bun_G \times X^I\]
    where the first map is a closed embedding and the second is ind-proper.
    \item[(ii)] There is an open embedding
      \[r_{G,I}^\pol : \HH_{G,I} \times \BB Z_G \mon \Bunc_{G,I}^\pol.\]
    \item[(iii)] The Hecke stack
      \[\HH_{G,I} \underset{X^I}{\times} \HH_{G,I} \to \Bun_G \times \Bun_G \times X^I\] 
    acts on $\Bunc_{G,I}^\pol$, and both action maps
      \[\HH_{G,I} \underset{\Bun_{G,I}}{\times} \Bunc_{G,I}^\pol \to \Bunc_{G,I}^\pol  \qquad  \Bunc_{G,I}^\pol \underset{\Bun_{G,I}}{\times} \HH_{G,I} \to \Bunc_{G,I}^\pol\]
    are ind-proper.
\end{itemize}
\end{proposition}

\subsubsection{} The verification of Proposition \ref{proposition:Hecke_comp_propoerties} is straightforward, but we will include the details in the next subsection. For now let us simply describe the fiber of $\Vin\Bun_{G,I}^\pol$ over a point $x \in X^I$ in the case $G = \SL_2$. It classifies a pair of principal $\SL_2$-bundles along with a nonzero rational map
    \[\alpha : \EE_1 \dashrightarrow \EE_2\]
of coherent sheaves which is regular away from $x$. 

\begin{proposition}
The functor
\[
    \Sph_{G,I} \to \DMod\left(\Bunc_{G,I}^\pol \right)
    \qquad \SS \mapsto r_{G,I,!}^\pol \cdot \HH(\SS)
\]
is canonically a map of $\Sph_{G,I}$-bimodules over $X^I$.
\end{proposition}

\subsubsection{} In what follows, we will mainly be interested in the full subcategory 
    \[\DMod\left(\Bunc_{G,I}^\pol \right)^\good \mon \DMod\left(\Bunc_{G,I}^\pol \right)\]
generated by the image of this functor.

\subsubsection{Proof} The functor in question is automatically colax monoidal, so it suffices to show that the colax morphisms
  \[r_{G,I,!}^\pol(\SS \star \id_{\Sph_{G,I}}) \to \SS \star r_{G,I,!}^\pol(\id_{\Sph_{G,I}})  \qquad  r_{G,I,!}^\pol(\id_{\Sph_{G,I}} \star \SS) \to r_{G,I,!}^\pol(\id_{\Sph_{G,I}}) \star \SS\]
are isomorphisms. Indeed, this would imply that every object of $\Sph_{G,I}$ is strict for both the left and right actions. This then implies that the functor is strictly linear.

\subsubsection{} By symmetry, we only need to consider the right action. Consider the diagram
\[
\begin{tikzcd}[
  column sep={4em,between origins},
  row sep={4em,between origins}
]
    & \HH_{G,I} \arrow[dl] \arrow[rr, "\aph r_G"] & & \Bunc_G \!\underset{\Bun_G}{\times}\! \HH_{G,I} \arrow[dl] \arrow[dr] & \\
    \Bun_G \arrow[rr, "r_G"] & & \Bunc_G \arrow[dr] & & \HH_{G,I} \arrow[dl] \\
     & & & \Bun_G &
\end{tikzcd}
\]
The colax constraint for the right action is the pushforward of
  \[{\aph r_G}_! \cdot \HH(\SS) \to r_{G,!} \shoximes \HH(\SS) \]
along the proper morphism
\[
    \Bunc_G \underset{\Bun_G}{\times} \HH_{G,I} 
    \simeq \Bunc_{G,I} \underset{\Bun_{G,I}}{\times} \HH_{G,I}
    \to \Bunc_{G,I}^\pol.
\]
This map is an isomorphism because $\HH(\SS)$ is ULA over $\Bun_G$.

\subsubsection{Notation} 
\begin{itemize}
    \item We write $(-)_\HH$ and $(-)^\trans$ for restriction to the open loci
\[
    \Bunc_{N,I}^\pol \underset{\Bun_{G,I}}{\times} \HH_{G,I}
    \mon \left(\Bunc_{N,I}^\pol \underset{\Bun_{G,I}}{\times} \Bunc_{G,I}^\pol \right)^\trans 
    \mon \Bunc_{N,I}^\pol \underset{\Bun_{G,I}}{\times} \Bunc_{G,I}^\pol.
\]
    \item We have the ind-representable morphism
    \[
    \pi_2 : \left(\Bunc_{N,I}^\pol \underset{\Bun_{G,I}}{\times} \Bunc_{G,I}^\pol \right)^\trans
    \to \Bun_{G,I}.
    \]
\end{itemize}

\subsubsection{} Now define a pair of functors
\[
    \Whit(G)_I^\glob \otimes \DMod\left(\Bunc_{G,I}^\pol \right)^\good
    \to \DMod(\Bun_{G,I}).
\]
by
\[
    \Poinc_{I,!}^\ext(\WW, \FF) = \pi_{2,!} \left( \WW \shoximes \FF \right)_\HH \qquad
    \Poinc_{I,?}^\ext(\WW, \FF) = \pi_{2,*} \left( \WW \shoximes \FF \right)^\trans
\]
By forgetting supports, we obtain a natural transformation
  \[\eta_{G,I}^\ext : \Poinc_{I,!}^\ext \to \Poinc_{I,*}^\ext.\]

\begin{proposition}
We have:
\begin{itemize}
    \item[(i)] The functor $\Poinc_{I,!}^\ext$ descends to a $\Sph_{G,I}$-linear functor
        \[\Whit(G)_I^\glob \underset{\Sph_{G,I}}{\otimes} \DMod\left(\Bunc_{G,I}^\pol \right)^\good \to \DMod(\Bun_{G,I})\]
    \item[(ii)] The functor $\Poinc_{I,*}^\ext$ descends to a $\Sph_{G,I}$-linear functor
      \[\Whit(G)_I^\glob \underset{\Sph_{G,I}}{\otimes} \DMod\left(\Bunc_{G,I}^\pol \right)^\good \to \DMod(\Bun_{G,I}).\]
    \item[(iii)] Furthermore, $\eta_{G,I}^\ext$ descends to a map of $\Sph_{G,I}$-linear functors.
\end{itemize}
\end{proposition}

\subsubsection{Proof of (i)} Note that $\Poinc_{I,!}^\ext$ admits a factorization
\[
    \Whit(G)_I^\glob \otimes \DMod\left( \Bunc_{G,I}^\pol \right)^\good
    \to \Whit(G)_I^\glob \otimes \DMod(\HH_{G,I})^\good
    \to \DMod(\Bun_{G,I})
\]
where 
    \[\DMod(\HH_{G,I})^\good \mon \DMod(\HH_{G,I})\]
denotes the full subcategory generated by the image of
    \[\HH(-) : \Sph_{G,I} \to \DMod(\HH_{G,I}).\]
The first functor in this factorization is linear for all three $\Sph_{G,I}$ actions, and the second functor factors through a $\Sph_{G,I}$-linear functor
\[
    \Whit(G)_I^\glob \underset{\Sph_{G,I}}{\otimes} \DMod(\HH_{G,I})^\good
    \to \DMod(\Bun_{G,I}).
\]

\subsubsection{Proof of (ii)} Actually the analogous claim for all of $\DMod\left(\Bunc_{G,I}^\pol \right)$ holds by base change.

\subsubsection{Proof of (iii)} This follows from the compatibility of forgetting supports with proper pushforwards.

\subsubsection{Proof of Proposition \ref{proposition:Sph_linear}} Form the map
\[
    \Whit(G)_I^\glob \simeq \Whit(G)_I^\glob \underset{\Sph_{G,I}}{\otimes} \Sph_{G,I} 
    \to \Whit(G)_I^\glob \underset{\Sph_{G,I}}{\otimes} \DMod\left(\Bunc_{G,I}^\pol \right)^\good
\]
of right $\Sph_{G,I}$-modules. Its composite with $\eta_{G,I}^\ext : \Poinc_{I,!}^\ext \to \Poinc_{I,*}^\ext$ identifies with $\eta_{G,I}$.

\subsection{Proof of Proposition \ref{proposition:Hecke_comp_propoerties}}

\subsubsection{Proof of (i)} To simplify notation, we will write the proof for a fixed point $x \in X^I$. Let $\UU = X \setminus x$. We claim that that both of the maps
    \[
    \Vin\Bun_G \to \Vin\Bun_G^\pol \to \Vin\Bun_G(\UU)
    \]
are closed immersions. This follows from the fact that the morphism
    \[\Maps(X, S/H) \to \Maps(\UU, S/H)\]
is a closed immersion for any finite scheme $S$ with a group action by $H$. Indeed, the second arrow is a base change of
    \[\Maps(X, G \bs T_\ad^+ / G) \to \Maps(\UU, G \bs T_\ad^+ / G)\]
while the composite is a base change of
    \[\Maps(X, G \bs \Vin_G / G) \to \Maps(\UU, G \bs \Vin_G / G).\]
The first arrow is then a closed immersion by the two out of three property. Now take quotients by $T$ to obtain the closed immersions
    \[\Bunc_G \mon \Bunc_G^\pol \mon \Bunc_G(\UU).\]
Replacing a complete curve with an open curve in the proof \cite[Appendix A]{VinGr} that $\Bunc_G$ is proper over $\Bun_G \times \Bun_G$ shows that $\Bunc_G(\UU)$ is ind-proper over $\Bun_G(\UU) \times \Bun_G(\UU)$. Hence $\Bunc_G^\pol$ is proper over the closed substack
    \[\Bun_G \times \Bun_G \mon \Bun_G(\UU) \times \Bun_G(\UU).\]

\subsubsection{Proof of (ii)} Note that $\HH_{G,I}$ identifies with the base change of
    \[\Vin\Bun_{G,I}^\pol \to \Maps(X, T_\ad^+) \simeq T_\ad^+\]
along the open subset $T_\ad \mon T_\ad^+$. Taking the quotient by $T$ gives $r_{G,I}^\pol$.

\subsubsection{Proof of (iii)} The existence of the action in (iii) is clear and the ind-properness follows from that of the affine Grassmannian.

\subsection{The Ran version of the Jacquet functor}

\subsubsection{} The purpose of this subsection is to construct the Ran version of the Jacquet functor along with some of its attendant compatibilities. Recall that we have the space
    \[\Ran_\Ran \mon \Ran \times \Ran\]
classifying pairs of points $x \subset y$ in $\Ran \times \Ran$. The Ran version of the Jacquet functor will be a map        \[\Jac_\Ran : \Whit(G)_\Ran \to \Whit(M)_{\Ran_\Ran},\]
where $\Whit(M)_{\Ran_\Ran}$ is the pullback of $\Whit(M)$ along the projection onto the second copy of $\Ran$. 

\subsubsection{} We will work over powers of the curve. As before, the definition of this functor will actually be as a composite of functors
    \[\Jac_I : \Whit(G)_I \to \Whit(M)_{\Conf_I} \to \Whit(M)_{\Ran_I}\]
Here is a complete list of the objects we need to construct in this subsection:
\begin{itemize}
    \item A Jacquet functor
        \[\Jac_{\Conf,I} : \Whit(G)_I \to \Whit(M)_{\Conf_I}\]
    with the compatibility
        \begin{equation}
        \label{eqn:Jac_Conf_I_Poinc}
        \begin{tikzcd}[
          column sep={8em,between origins},
          row sep={4em,between origins}
        ]
            \Whit(G)_I \arrow[r, "\Jac"] \arrow[d, "\Poinc_!"] & \Whit(M)_{\Conf_I} \arrow[d, "\Poinc_!"] \\
            \DMod(\Bun_G)_I \arrow[r, "\CT_!^-"] & \DMod(\Bun_M)_I
        \end{tikzcd}
        \end{equation}
    \item A co-Jacquet functor    
        \[\co\Jac_{\Conf,I} : \Whit(M)_{\Conf_I} \to \Whit(G)_I\]
    with the compatibility
        \begin{equation}
        \label{eqn:Jac_co_Conf_I_coeff}
        \begin{tikzcd}[
          column sep={8em,between origins},
          row sep={4em,between origins}
        ]
            \DMod(\Bun_M)_I \arrow[r, "\Eis_!^-"] \arrow[d, "\coeff_*"] & \DMod(\Bun_G)_I \arrow[d, "\coeff_*"] \\
            \Whit(M)_{\Conf_I} \arrow[r, "\Jac_\co"] & \Whit(G)_I
        \end{tikzcd}
        \end{equation}        
    \item A duality
        \[\DD : \Whit_\psi(M)_{\Conf_I} \simto \Whit_{-\psi}(M)_{\Conf_I}^\vee\]
    with the compatibility
        \begin{equation}
        \label{eqn:Jac_I_duality}
        \begin{tikzcd}[
          column sep={10em,between origins},
          row sep={4em,between origins}
        ]
            \Whit_\psi(G)_I \arrow[r, "\Jac"] \arrow[d, "\DD{[2\delta_G]}" left] & \Whit_\psi(M)_{\Conf_I} \arrow[d, "\DD{[2\delta_M]}"] \\
            \Whit_{-\psi}(G)_I^\vee \arrow[r, "\Jac_\co^\vee"] & \Whit_{-\psi}(G)_{\Conf_I}^\vee
        \end{tikzcd}
        \end{equation}
    \item A functor
        \[\sprd_{I,*} : \Whit(M)_{\Conf_I} \to \Whit(M)_{\Ran_I}\]
    with the compatibility
        \begin{equation}
        \label{eqn:sprd_I_duality}
        \begin{tikzcd}[
          column sep={11em,between origins},
          row sep={4em,between origins}
        ]
            \Whit_\psi(M)_{\Conf_I} \arrow[r, "\sprd_*"] \arrow[d, "\DD" left] & \Whit_\psi(M)_{\Ran_I} \arrow[d, "\DD"] \\
            \Whit_{-\psi}(M)_{\Conf_I}^\vee \arrow[r, "(\sprd^!)^\vee"] & \Whit_{-\psi}(M)_{\Ran_I}^\vee
        \end{tikzcd}
        \end{equation}
    and an identification
        \begin{equation}
        \label{eqn:sprd_I_Poinc}
            \Poinc_{\Conf_I,!} \simeq \Poinc_{\Ran_I,!} \cdot \sprd_*
        \end{equation}
    of functors
        \[\Whit(M)_{\Conf_I} \to \Whit(M)_{\Ran_I}.\]
    Here $\sprd_I^!$ is the right adjoint to $\sprd_{I,*}$.
\end{itemize}

\subsubsection{Pole-adding} Let $\YY \to X^I$ be any prestack equipped with an action of the constant monoid scheme $(\Lambda^\pos)^I$. We will think of and refer to such actions as \emph{zero-adding}. Inverting this action gives a new prestack denoted $\YY^\pol$. All of our examples will have the feature that the action maps
  \[\lambda : \YY \to \YY  \qquad \lambda \in (\Lambda_{G,P}^\pos)^I.\]
are closed immersions.

\subsubsection{} Let us list a number of actions to which we will apply the pole-adding procedure. In many cases, the action will factor through the map $\Lambda^\pos \epi \Lambda_{G,P}^\pos$.
\begin{itemize}
    \item We give $\Bun_{G,I}$ and $\Bunt_{P^-, I}$ the trivial actions. So we will not write $\Bun_{G,I}^\pol$ or $\Bunt_{P^-, I}^\pol$.
    \item We give $\Conf_{G,P} \times X^I$ the $(\Lambda_{G,P}^\pos)^I$-action of
        \[\lambda \cdot (D, x) = (D + \lambda \cdot x, x).\]
    To reduce notation, we write
        \[\Conf_I = (\Conf_{G,P} \times X^I)^\pol.\]
    \item We give $\Bunc_{N,I}$ the usual action of $(\Lambda^\pos)^I$ by adding zeros. Note that the map $\Bunc_{N,I} \to \Bun_{G,I}$ is equivariant for this action. We warn that this action is not compatible with the formation of Whittaker categories.
    \item In the same way,
        \[\Bunc_{N_M, \Conf}^\pol \times X^I\]
    has an action of $\Lambda_{G,P}^\pos$ by adding zeros for which the projection to $\Conf \times X^I$ is equivariant. We write
        \[\Bunc_{N_M, \Conf_I}^\pol \to \Conf_I\]
    for the result of applying pole-adding to this map.
    \item The open substack
        \[
        \ZZ_{\Nc, \Pt^-, I}
        \mon \left(\Bunc_N \underset{\Bun_G}{\times} \Bunt_{P^-}\right)_I
        \]
    is stable under zero-adding, and the map
        \[
        \ZZ_{\Nc, \Pt^-, I}
        \to \Bunc_{N_M, \Conf}^\pol \times X^I
        \]
    is equivariant for this action.
\end{itemize}

\subsubsection{} The pole-adding procedure then gives a polar variant of (\ref{eqn:zastava_diagram}).
\[
\begin{tikzcd}[
  column sep={4em,between origins},
  row sep={4em,between origins}
]
     & \ZZ_{\Nc, \Pt^-, I}^\pol \arrow[dl, "v" above left] \arrow[dr, "t"] \arrow[drrr, dashed, bend left=15, "u"] & &  & \\
    \Bunc_{N,I}^\pol \arrow[dr] & & \Bunt_{P^-, I} \arrow[dl] \arrow[dr] & & \Bunc_{N_M, \Conf_I}^\pol \arrow[dl] \\
     & \Bun_{G,I} & & \Bun_{M,I} &
\end{tikzcd}
\]
We define
\[
    \Jac_\Conf : \Whit(G)_I \to \Whit(M)_{\Conf_I}
    \qquad \WW \mapsto u_* \cdot \left( \WW \shoximes j_{P,!}^- \right)_\ZZt [2 \delta_G]
\]
where the formation of $\shoximes$ is with respect to the fiber square
    \[
    \begin{tikzcd}[
  column sep={3em,between origins},
  row sep={3em,between origins}
]
        & \square \arrow[dl] \arrow[dr] \\
      \Bunc_{N,I} \arrow[dr] & & \Bunt_{P^-} \arrow[dl] \\
       & \Bun_G
    \end{tikzcd}
    \]
The fact that this formula produces an object in the Whittaker category is checked in the same way as Proposition \ref{proposition:Jac_Whit_preservation}. Now we have as before the natural transformations
\begin{align*}
    \CT_{I,!}^- \cdot \Poinc_{G,I,!}
    & \overset{\alpha}{\longleftarrow} \pi_! \cdot u_! \left( - \stoximes j_{P,!}^- \right)_\ZZt \\
    & \qquad\qquad \overset{\beta}{\longrightarrow} \pi_! \cdot u_* \left( - \stoximes j_{P,!}^- \right)_\ZZt 
    \overset{\gamma}{\longrightarrow} \Poinc_{M,I,!} \cdot \Jac_{\Conf,I}.
\end{align*}

\begin{proposition}
All of these maps are isomorphisms.
\end{proposition}

\subsubsection{Proof} The vanishing argument used to prove Proposition \ref{proposition:zast_vanishing} shows that $\alpha$ is an isomorphism. Also $\beta$ is an isomorphism because $u$ is proper: it was obtained from a proper map via the pole-adding procedure. We will prove that $\gamma$ is an isomorphism by use of the following reductions.

\subsubsection{First reduction: to a point} The four functors appearing above are colax linear over $\DMod(X^I)$. We claim that they are actually strictly linear. For $\CT_{I,!}^- \cdot \Poinc_{G,I,!}$, this follows from the second adjointness $\CT_{I,!}^- \simeq \CT_{I,*}$ and the fact that $\Poinc_{G,I,!}$ is even $\Sph_{G,I}$-linear. Since $\alpha$ and $\beta$ are isomorphisms, it remains to observe $\Poinc_{M,I,!} \cdot \Jac_\Conf$ is linear because both of its factors are. Therefore, to check that $\gamma$ is an equivalence, we may do so stratawise.

\subsubsection{Second reduction: to $\WW_{\vac, x}$} The restriction of $\Whit(G)_I$ to each stratum of $X^I$ is actually constant, so it suffices to check that $\gamma$ is an equivalence pointwise. For this it is enough to show that the purity morphism
    \[
    \left( - \stoximes j_{P,!}^- \right)_\ZZt 
    \to \left( - \shoximes j_{P,!}^- \right)_\ZZt \, [2 \delta_{G}]
    \]
defines an isomorphism of functors
    \[
    \Whit(G)_x \to \DMod\left(\ZZ_{\Nc,\Pt,x}^\pol \right).
    \]
The first functor is $\Sph_{G,x}$-linear by the ULA property of Hecke operators and the second is obviously so. Therefore, the result follows from Lemma \ref{lem:zast_ULA}, which says that the purity morphism is an isomorphism on the object $\WW_{\vac, x}$.

\subsubsection{Dualities} 

As before, we have a $\DMod(X^I)$-linear perfect pairing
\[
    \Whit_\psi(M)_{\Conf_I}^\glob \otimes \Whit_{-\psi}(M)_{\Conf_I}^\glob \to \DMod(X^I)
\]
given by $!$-tensoring the underlying sheaves followed by $*$-pushforward along the map
  \[\Bunc_{N, \Conf_I}^\pol \to \Conf_I \to X^I.\]
We take $\Jac_\co$ to be the descent
\[
\begin{tikzcd}[
  column sep={9em,between origins},
  row sep={4em,between origins}
]
    \DMod\left( \Bunc_{N,I}^\pol \right) \arrow[r] \arrow[d] & \DMod\left( \Bunc_{N_M, \Conf_I}^\pol \right) \arrow[d] \\
    \Whit(G)_I \arrow[r, dashed, "\Jac_\co"] & \Whit(M)_{\Conf_I} 
\end{tikzcd}
\]
where the upper arrow is
  \[\WW \mapsto v_* \left( t^! \cdot j_{P,!}^- \shimes u^! \cdot \WW \right)[2 \delta_{M}].\]
It is exchanged with $\Jac$ under the above duality.

\subsubsection{Conf v Ran} We need a version of the diagram (\ref{eqn:Conf_Ran_vac}) with marked points.
\[
\begin{tikzcd}[
  column sep={8em,between origins},
  row sep={2.5em,between origins}
]
    & \Bunc_{N_M, \Ran_I}^{\neg, \pol} \arrow[dl, "\aph\pi"] \arrow[dr, "\aph i^\neg"] \arrow[dd] & \\
  \Bunc_{N_M, \Conf_I}^{\pol} \arrow[dd] & & \Bunc_{N_M, \Ran_I}^{\pol} \arrow[dd] \\
    & \Gr_{M^\ab, \Ran_I}^{\omega,\neg, \pol} \arrow[dl] \arrow[dr] &  \\
  \Conf_I & & \Gr_{M^\ab, \Ran_I}^\omega
\end{tikzcd}
\]
This diagram is constructed by applying the pole-adding procedure to the diagram
\[
\begin{tikzcd}[
  column sep={8em,between origins},
  row sep={2.5em,between origins}
]
    & \square \arrow[dl] \arrow[dr] \arrow[dd] & \\
  \Bunc_{N_M, \Conf}^\pol \times X^I \arrow[dd] & & \Bunc_{N_M, \Ran_I}^{\pol} \arrow[dd] \\
    & \Gr_{M^\ab, \Ran_I}^{\omega,\neg} \arrow[dl, "\pi" above left] \arrow[dr] &  \\
  \Conf \times X^I & & \Gr^\omega_{M^\ab, \Ran_I}
\end{tikzcd}
\]
Here are the definitions:
\begin{itemize}
    \item The right face of this diagram is the base change of the right face of (\ref{eqn:Conf_Ran_vac}) along the map $\Ran_I \to \Ran$. The right vertical arrow carries an action of the group scheme $\Gr_{M^\ab, X^I}$, and we let $(\Lambda_{G,P}^\pos)^I$ act on this arrow through the map
      \[(\Lambda_{G,P}^\pos)^I \times X^I \to \Gr_{M^\ab, X^I}\]
    of monoids over $X^I$. This $(\Lambda_{G,P}^\pos)^I$-action preserves the closed sub-prestack $\Gr_{M^\ab, \Ran_I}^{\omega,\neg}$. This completes the definition of the right face and shows that the adding poles procedure does not affect the right vertical arrow.
    \item The leftward map $\pi$ is given by the projections
      \[
      \Gr_{M^\ab, \Ran_I}^{\omega,\neg} \to \Gr_{M^\ab, \Ran}^{\omega,\neg} \to \Conf
      \qquad
      \Gr_{M^\ab, \Ran_I}^\neg \to \Ran_I \to X^I.
      \]
    As before, the fiber products of the two faces are canonically identified.
\end{itemize}
Now we define
    \[
    \sprd_{I,*} : \Whit(M)_{\Conf_I} \to \Whit(M)_{\Ran_I}
    \qquad \WW \mapsto (\aph i^\neg)_* \cdot \aph\pi^! \cdot \WW.
    \]
One obtains the compatibilities (\ref{eqn:sprd_I_duality}) and (\ref{eqn:sprd_I_Poinc}) exactly as in \S\ref{subsection:sprd_compatibilities}.